% Generated draft: requires mathematical and source review.
\documentclass[11pt]{article}
\usepackage[margin=25mm]{geometry}
\usepackage{fontspec}
\setmainfont{Latin Modern Roman}
\usepackage{amsmath,amssymb,mathtools,unicode-math}
\setmathfont{Latin Modern Math}
\usepackage{enumitem,needspace,xurl,hyperref}
\setlength{\parindent}{0pt}
\setlength{\parskip}{5pt}
\setlength{\emergencystretch}{3em}
\newcommand{\N}{\mathbb{N}}
\newcommand{\Z}{\mathbb{Z}}
\newcommand{\Q}{\mathbb{Q}}
\newcommand{\R}{\mathbb{R}}
\newcommand{\C}{\mathbb{C}}
\newcommand{\F}{\mathbb{F}}
\newcommand{\A}{\mathbb{A}}
\newcommand{\PP}{\mathbb{P}}
\newcommand{\E}{\mathbb{E}}
\newcommand{\eps}{\varepsilon}
\newcommand{\dd}{\,\mathrm{d}}
\begin{document}
\section*{UnsolvedMath ID: 30003314}
\textbf{Source problem code:} catalogue entry
\subsection{Definitions and mathematical statement}
Let \(\mathbf{No}\) denote the ordered field of surreal numbers, regarded as a proper class, and let \(\Q\subset\mathbf{No}\) denote its canonical ordered subfield of rational numbers. A surreal number \(x\) has a sign-sequence representation \(s_x\colon\alpha_x\to\{-,+\}\), where \(\alpha_x\) is an ordinal. The standard tree simplicity relation is given by \(x\prec y\) when \(s_x\) is a proper initial segment of \(s_y\). The supplied record does not specify its precise convention for simplicity or the intended extension of that convention to functions.

The original statement contains the following six questions. Its integration and integer-part questions are absent from the shorter cleaned statement and are retained here.
\begin{enumerate}
\item \textbf{Extension of a degree map to a valuation.} For an ordered abelian group \(\Gamma\), written additively, a field valuation has the form
\[
 v\colon\mathbf{No}\longrightarrow\Gamma\cup\{\infty\},
\]
with \(\infty>\gamma\) and \(\infty+\gamma=\infty\) for every \(\gamma\in\Gamma\), and satisfies, for all \(x,y\in\mathbf{No}\),
\[
 v(x)=\infty\iff x=0,\qquad v(1)=0,
\]
\[
 v(xy)=v(x)+v(y),\qquad
 v(x+y)\geq\min\{v(x),v(y)\}.
\]
The question asks whether the source's map \(\deg\) extends to such a valuation. The record supplies neither the definition and original domain of \(\deg\) nor its codomain, normalization, or compatibility with this valuation convention. Consequently, an exact extension condition cannot yet be stated. The permitted size of the value group and the foundational interpretation of a valuation on the proper class \(\mathbf{No}\) also remain unspecified.

\item \textbf{Integration and measure.} Can integration and measure theory on \(\mathbf{No}\) be developed beyond the semialgebraic and globally subanalytic categories? This is a research-program question as supplied: the record specifies no larger class of sets or functions, measurable structure, codomain of a measure or integral, additivity requirements, or desired integration axioms. Thus it does not determine a single existence proposition.

\item \textbf{Rational sign sequences.} Find a characterization, expressed in terms of the sign sequence \(s_x\), equivalent to
\[
 x\in\Q\subset\mathbf{No}.
\]
The record does not specify what form or effectiveness requirements such a characterization must have; a mere restatement using the represented surreal number would not identify the intended target.

\item \textbf{Simplicity for functions.} Can the simplicity order on surreal numbers be extended to functions? The domain and codomain of the functions, the class of functions under consideration, the order axioms, and the intended meaning of extension are not supplied. Hence no particular function-order construction or existence assertion is fixed by the record.

\item \textbf{Simplicity of the exponential.} Is the surreal exponential \(\exp_{\mathbf{No}}\) the simplest suitable homomorphism? The homomorphism aspect can be expressed as a map between the additive and positive multiplicative groups,
\[
 E\colon(\mathbf{No},+)\longrightarrow(\mathbf{No}_{>0},\cdot),
 \qquad E(x+y)=E(x)E(y).
\]
However, the record does not list the additional required properties, identify the competing homomorphisms, or define the function simplicity relation. It therefore does not specify a precise minimality or uniqueness statement for \(\exp_{\mathbf{No}}\).

\item \textbf{An explicit integer part satisfying true arithmetic.} An integer part of \(\mathbf{No}\) is a discretely ordered subring \(I\subseteq\mathbf{No}\), containing \(1\) as its least positive element, such that
\[
 \forall x\in\mathbf{No}\ \exists i\in I\quad i\leq x<i+1.
\]
With \(I_{\geq0}=\{i\in I:i\geq0\}\), the usual first-order interpretation of the arithmetic requirement is
\[
 (I_{\geq0};0,1,+,\cdot,<)
 \models
 \operatorname{Th}(\N;0,1,+,\cdot,<).
\]
The original question asks for an explicit construction of such an integer part. The supplied background additionally describes requests for explicit constructions without choice; that qualification must be retained, although the record does not specify its allocation among the component questions or its foundational meaning. In particular, the permitted choice principles, class framework, and criterion for an explicit construction are unresolved.
\end{enumerate}

These are statement-only formulations and explicit records of missing specifications. The supplied open label and literature assessment do not establish that every component question is open, and the general existence of integer parts mentioned in the supplied assessment does not itself settle the explicit-construction and arithmetic requirements above.
\subsection{Short English statement}
The source collects six questions about surreal numbers: extending a degree map to a valuation, developing integration and measure beyond the stated categories, recognizing rational numbers from sign sequences, extending simplicity to functions, characterizing the exponential by simplicity, and explicitly constructing an integer part satisfying true arithmetic. Several essential definitions and requirements are missing. The shorter cleaned statement omits two questions that are preserved here from the original. The supplied background also raises construction without choice, without specifying its exact scope.
\subsection{Sources}
\begin{itemize}
\item Supplied local UnsolvedMath registry record, permanent ID \(30003314\), title \emph{Valuation and Simplicity Problems for the Surreal Numbers}, source problem code \texttt{catalogue entry}, in the Oberwolfach Reports Open Problems collection. No individual proposer is supplied; the proposed year is recorded as 2016.
\item Supplied report citation: \emph{Mini-Workshop: Surreal Numbers, Surreal Analysis, Hahn Fields and Derivations} (2017), Oberwolfach Reports, pp.\ 3313--3372, DOI \texttt{10.4171/owr/2016/60}. Original OWR report: \url{https://doi.org/10.4171/owr/2016/60}. The record describes the report as documenting workshop discussions and open problems. The relevant passage and its individual attribution are not reproduced in the supplied record.
\item Supplied literature reference, labeled Carl--Krapp: \url{https://doi.org/10.4115/jla.2021.13.3}. No title, passage, or precise supporting result is supplied.
\item Supplied literature reference, labeled Mourgues--Ressayre: \url{https://doi.org/10.2307/2275224}. No title, passage, or precise supporting result is supplied.
\item Literature assessment dated \texttt{2026-08-21}: the supplied status is open. The supplied assessment says that some component questions have positive results, including existence of integer parts for real closed fields, while combined requests concerning \(\mathbf{No}\), explicit constructions without choice, and remaining operations are unresolved. This is supplied dated curation, not independently verified literature assessment; the source itself requests independent verification of time-sensitive status and cited claims.
\item The registry labels the statement \texttt{corrected\_verified} and reports that an official report passage was checked and rewritten as a self-contained problem while preserving mathematical content. This is a supplied verification claim, not verification performed for this draft. The discrepancy between the shorter statement and the original six-part statement remains material.
\end{itemize}
\end{document}
